\documentclass[11pt,a4paper]{article}

\usepackage{fontspec}
\usepackage{polyglossia}
\setmainlanguage{russian}
\setotherlanguage{english}
\setmainfont{Times New Roman}
\setsansfont{Arial}
\setmonofont{Courier New}
\newfontfamily\cyrillicfont{Times New Roman}
\newfontfamily\cyrillicfontsf{Arial}
\newfontfamily\cyrillicfonttt{Courier New}
\usepackage[a4paper,margin=2.4cm]{geometry}
\usepackage{microtype}
\usepackage{amsmath,amssymb}
\usepackage[hidelinks]{hyperref}
\usepackage{setspace}
\usepackage{enumitem}
\usepackage{titlesec}
\usepackage{graphicx}
\usepackage[font=small,labelfont=bf]{caption}
\usepackage{xcolor}
\usepackage{longtable,booktabs,array}
\newcolumntype{L}[1]{>{\raggedright\arraybackslash}p{#1}}
\usepackage{float}
\usepackage{placeins}
\usepackage{tcolorbox}
\tcbuselibrary{skins,breakable}
\usepackage{tikz}
\usetikzlibrary{positioning,arrows.meta}

\setstretch{1.08}
\setlength{\parskip}{0.45em}
\setlength{\parindent}{0pt}

\titleformat{\section}{\large\bfseries}{\thesection.}{0.6em}{}
\titleformat{\subsection}{\normalsize\bfseries}{\thesubsection.}{0.5em}{}

\hypersetup{
  pdftitle={Euclidean Condensate Theory — Russian Summary},
  pdfauthor={Valeriy Blagovidov},
  pdflang={ru}
}

\newtcolorbox{mechbox}[1][] {
  enhanced, breakable,
  colback=gray!4, colframe=gray!55,
  boxrule=0.4pt, arc=2pt,
  left=8pt, right=8pt, top=4pt, bottom=4pt,
  fonttitle=\bfseries,
  before skip=6pt, after skip=8pt,
  #1
}

\newcommand{\statA}{\textbf{Уровень~A}}
\newcommand{\statB}{\textbf{Уровень~B}}
\newcommand{\statC}{\textbf{Уровень~C}}
\newcommand{\statO}{\textbf{Открыто}}

\title{\textbf{Euclidean Condensate Theory (ECT):\\
Краткий обзор механизмов и ключевых результатов}}

\author{Valeriy Blagovidov\thanks{Для переписки:
\texttt{blagovidov@phystech.edu},
\texttt{vblagovidov@gmail.com}.\,
ORCID: \href{https://orcid.org/0009-0008-6707-7068}{0009-0008-6707-7068}.}}
\date{Июнь 2026}

\begin{document}
\maketitle

\begin{abstract}
\noindent
Euclidean Condensate Theory (ECT) исследует возможность того,
что лоренцево пространство-время, когерентная квантовая
динамика, эффективный гравитационный сектор и отдельные
галактические и космологические феномены возникают как
секторно-зависимые описания одной упорядоченной фазы: реальной
скалярной среды $\Phi$ на четырехмерном евклидовом носителе со
спонтанным нарушением симметрии $O(4)\!\to\!O(3)$.
Центральное утверждение носит архитектурный характер: структуры,
которые обычно постулируются по отдельности, могут оказаться
разными режимами одной нарушенной ветви.
Этот документ задуман как краткая точка входа для физика,
которому нужны основная конструкция, ключевые результаты,
статусный уровень каждой претензии (Уровень~A/B/C/Открыто) и
главные проверочные точки без немедленного обращения к полному
препринту.
Цветовой сектор рассматривается отдельно: точная локальная
симметрия $SU(3)$ требует дополнительного постулата завершения
P7 и не входит в минимальный $O(4)$-упорядоченный конденсат.
Полный препринт доступен по адресу
\href{https://doi.org/10.5281/zenodo.18917929}{doi:10.5281/zenodo.18917929},
а сопроводительный обзор — по адресу
\href{https://doi.org/10.5281/zenodo.19430795}{doi:10.5281/zenodo.19430795}.
\end{abstract}

% =====================================================================
\section{Введение: от евклидовой среды к наблюдаемой физике}
\label{sec:intro}
% =====================================================================

Цель Euclidean Condensate Theory (ECT) — проверить, могут ли
несколько структур, которые сегодня вводятся в фундаментальной
физике по отдельности, возникать из одной и той же
упорядоченной фазы. Время, причинный конус, лоренцева метрика,
квантовое описание когерентной материи и гравитационный сектор
обычно наслаиваются друг на друга через независимые постулаты.
ECT предлагает рассматривать их вместо этого как эффективное
описание одной нарушенной фазы более примитивной структуры,
которая сама по себе не содержит ни времени, ни метрики, ни
квантового языка.

Кандидат на такую примитивную структуру намеренно минимален.
Это четырехмерный евклидов носитель, на котором задана одна
реальная скалярная среда $\Phi$. На этом уровне среда
допредфизична: в ней нет понятия энергии, постоянной Планка или
выделенного направления. Короткий список структурных свойств
задает, что именно допускает эта среда: связность, направленную
изотропию относительно четырехмерных вращений, локальную
регулярность при допустимых дефектных подмногообразиях,
существование упорядоченных конфигураций и тот факт, что
внутренние координаты упорядоченной ветви являются
описательными избыточностями, а не физическими наблюдаемыми.

Центральный структурный шаг ECT — появление упорядоченного
основного состояния, в котором градиент $\Phi$ приобретает
ненулевое среднее значение вдоль одного из четырех евклидовых
направлений. Полная четырехмерная вращательная симметрия
нарушается до своей трехмерной подгруппы; одно направление
становится выделенным. Аналогия с фазами упорядочения в
сверхтекучем гелии-3A и с градиентно-конденсатными фазами в
системах типа бозе-эйнштейновского конденсата носит
методологический характер, а не характер отождествления: ECT
помещает этот механизм на фундаментальный уровень, а не
рассматривает его как материально-физический эффект, возникающий
поверх уже существующего пространства-времени.

Нарушенная фаза несет две качественно различные семьи
возмущений. Гладкие однозначные возмущения коллективной фазы
образуют \emph{когерентный сектор}: для него возможна частичная
реконструкция стандартного квантового формализма через
топологическую шкалу действия и мост отражательной
положительности, а в стационарном когерентном режиме
восстанавливается эволюция типа Шредингера. Медленные возмущения упорядоченного фона мотивируют
геометрическое/ориентационное расширение.  Контролируемое уравнение
здесь является скалярным.  Физический ориентационный стресс
$T^{(n)}_{\mu\nu}$ требует отдельно заданного ковариантного
$S_n[g,n,\lambda_n]$; нынешние анизотропные тензоры являются внешними
прокси, а ориентационный вклад остается \statO{}. Галактическое приложение использует отдельно заданную
феноменологическую интерполяцию Level~C; она не выводится из
скалярного действия или предварительного однопетлевого результата.
Космологические числа также принадлежат отдельному диагностическому
слою, а не вычисленному ориентационному вкладу.

В этой картине лоренцево пространство-время, квантовая
когерентность и гравитация выступают как
секторно-зависимые режимы одной упорядоченной среды, а не как
несвязанные исходные аксиомы.
ECT организована как общая архитектура упорядоченной среды с
отдельно отслеживаемым уровнем строгости для каждого сектора.
Цветовой калибровочный сектор лежит вне минимальной
четырехмерной упорядоченной структуры и вводится позже как
отдельный слой внутреннего завершения.

Этот обзор должен работать как рабочая карта механизмов,
утверждений, статусных уровней и фальсификаторов.
Его задача — дать читателю возможность быстро увидеть, какие
части конструкции являются структурными выводами, какие опираются
на предпосылки согласования или замыкания, а какие остаются
открытыми.
Четырехуровневая дисциплина, введенная в
\S\ref{sec:status_discipline}, сохраняет это различие явным во
всем документе.

\paragraph{Как читать этот обзор.}
Читателю, которого интересует концептуальная архитектура, стоит
начать с разделов \S\S\ref{sec:intro}--\ref{sec:logic}.
Читателю, которого прежде всего интересует тестируемая
феноменология, лучше начать с
\S\S\ref{sec:cosmology}--\ref{sec:falsifiers}.
Читателю, который хочет сразу увидеть возможные точки отказа,
следует перейти к \S\S\ref{sec:falsifiers}--\ref{sec:status_map}.
Полный препринт и сопроводительный обзор
(ссылки приведены в \S\ref{sec:resources}) содержат технические
выводы, лежащие в основе каждого из утверждений этого обзора.

% =====================================================================
\section{Объем конструкции: утверждения, открытые пункты и текущие границы}
\label{sec:scope}
% =====================================================================

\textbf{Семь структурных утверждений ECT}, сформулированные на
уровне механизма, а не формулы:

\begin{enumerate}[leftmargin=1.4em,itemsep=2pt]
\item Лоренцево пространство-время и единый универсальный
      низкоэнергетический причинный конус, эффективная скорость
      которого определяется анизотропной жесткостью
      упорядоченной ветви
      (\statA{} для возникновения конуса внутри нарушенной фазы;
      \statB{} для согласования скорости возникшего конуса с
      наблюдаемой вакуумной скоростью света).
\item Когерентные квантовые сектора с доквантовой топологической
      шкалой действия, отождествляемой с постоянной Планка на
      уровне согласования с наблюдаемой физикой (\statB), и
      частичная реконструкция
      стандартного квантового формализма через отражательную
      положительность и мост Остервальдера--Шрадера; связь
      спин--статистика имеет статус \statA{} для свободного
      когерентного ядра по тому же евклидову маршруту.
\item Контролируемое скалярное метрическое уравнение внутри
      объявленного анзаца с отдельным согласованием ньютоновской
      нормировки (\statB{} внутри скалярного анзаца).  Физический
      ориентационный стресс требует $S_n[g,n,\lambda_n]$ и остается
      \statO{}; $\mathcal X^{\rm cand}_{\mu\nu}$ — только внешний прокси.
\item \statA{}-структурное развязывание \emph{прямого}
      классического канала фоновой вакуумной энергии в
      космологической постоянной через эмерджентный унимодулярный
      механизм на упорядоченной ветви. Это структурное устранение
      катастрофы прямого канала, \emph{а не} вывод наблюдаемого
      позднего значения $\rho_\Lambda$.
\item Галактическая феноменология задается интерполяционным
      fit-замыканием (\statC{}); наклон BTFR~4 является условной
      алгеброй B/C внутри принятого функционала.  Скалярно-галактическое
      согласование, закон среды $\Xi$ и интерфейс UDG/cluster
      остаются \statO{}
      (\S\ref{sec:galactic}).
\item Сохраняемое пятипробное окно космологической совместимости
      для одного диагностического параметра деформации на уровне
      $1\sigma$ (\statC{}; оно явно подается как
      многоканальный диагностический тест согласованности, а не
      как предсказание из первых принципов).
\item Отдельно постулированный слой цветового завершения $SU(3)$
      (постулат~P7), введенный потому, что точная локальная
      симметрия $SU(3)$ не может быть получена из минимального
      четырехмерного упорядоченного конденсата.
\end{enumerate}

\textbf{Ключевые открытые пункты текущей версии ECT}
(отслеживаются как \statO{}, если не оговорено иное):

\begin{itemize}[leftmargin=1.4em,itemsep=2pt]
\item Строгая теорема из первых принципов для правила Борна и
      стороны проблемы измерения, связанной с единственным
      исходом. В ECT есть условный путь \statB{} к весам типа
      Борна через маршрут отражательной положительности и
      реконструкции Остервальдера--Шрадера (RP/OS), а также
      аргумент типа Глисона, если принять три предпосылки
      редуцированного
      описания: C1 — декогерированные альтернативы несут
      эффективную гильбертову / проекторную структуру; C2 — эти
      альтернативы приблизительно ортогональны; C3 — вероятности
      аддитивны по взаимоисключающим альтернативам.
      Вывод самих C1--C3 из голых постулатов ECT и проблема
      единственного исхода / правила обновления остаются \statO{}
      (см. \S\ref{sec:quantum}).
\item Корреляции белловского типа, условие отсутствия сигнализации
      (no-signalling) и предел
      Цирельсона не выведены из текущей структуры ECT;
      common-medium nonfactorisation является кандидатным путём
      \statB{}/\statO{}; физический state/channel, Bell/CHSH,
      no-signalling и предел Цирельсона остаются \statO.
\item Полная низкоэнергетическая квантово-полевая архитектура —
      строгая локальность, кластерное разложение, аналитичность,
      симметрия кроссинга и S-матричная теория — не выведены
      заново из
      формулировки через упорядоченную среду; согласование
      выполнено только на уровне эффективного действия.
\item Вывод из первых принципов наблюдаемого позднего значения
      космологической постоянной. ECT устраняет прямой канал
      фоновой вакуумной энергии, но само наблюдаемое значение
      остается отдельной проблемой уровня \statC/\statO{}.
\item Динамический вывод самой упорядоченной ветви из голого
      действия (градиентная конденсация, OP-grad). В текущем
      препринте эта программа разложена на три зафиксированные
      подзадачи, а оставшийся разрыв локализован в одном мосте
      релаксации ограничений / коарс-грейнинга; рассмотренные к
      настоящему моменту кандидаты на такой мост этот разрыв не
      закрывают, поэтому шаг остается \statO.
\item Вывод назначений гиперзаряда, полного наблюдаемого спектра
      электрических зарядов, структуры ароматов, матриц смешивания
      CKM/PMNS и иерархии трех поколений Стандартной модели.
\item Решение strong-CP проблемы: топологическая связь
      $\Theta(\phi)F\tilde F$ в ECT эту проблему
      \emph{возвращает}, а не устраняет.
\item Предсказательная КХД-подобная феноменология из цветового
      завершения: при постоянных фоновых константах связи слой P7
      остается структурно классифицирующим, а не предсказательным.
\item Предсказательный закон среды, замыкающий галактический
      интерфейс в стресс-тестах UDG, dwarf и bullet-cluster
      сопряженных пар (центральный операциональный пробел).
\item Завершенное ультрафиолетовое завершение всех секторов.
\end{itemize}

ECT поэтому лучше читать как структурированную теоретическую
рамку с определенным ядром, явными слоями согласования и конечным
списком открытых зависимостей.
Каждый положительный результат ниже снабжен явным статусным
маркером и указателем на соответствующее место в полном препринте.

% =====================================================================
\section{Архитектура модели}
\label{sec:architecture}
% =====================================================================

ECT построена на трех структурных слоях.

\textbf{S-слой (онтология $\Phi$-среды).}
Десять свойств S1--S10 описывают $\Phi$-среду в чисто
геометрических и конфигурационных терминах, до всякого
упорядочения, энергии или квантового языка:
связный четырехмерный евклидов носитель (S1);
$O(4)$-направленная изотропия (S2);
точечная однородность (S3);
локальная регулярность с допустимыми дефектными
подмногообразиями (S4);
локальность (S5);
множественность допустимых конфигураций (S6);
допустимость упорядоченных состояний с
$\langle\partial_A\Phi\rangle\neq 0$ (S7);
средовой характер с поточечными внутренними координатами как
описательными избыточностями (S8);
неоднородность упорядочения (S9);
топологическая полнота (S10).
Поскольку $\Phi$-среда является допредфизичной, на этом уровне
не используется ни язык энергии, ни язык квантовых флуктуаций.

\textbf{Вариационный принцип (DP).}
Физически реализуемые конфигурации экстремизируют локальный
евклидов функционал $S_E[\Phi]$, совместимый с базовыми
симметриями среды, при условиях, обеспечивающих корректную
постановку задачи.

\textbf{Рабочие постулаты P1--P6.}
Минимальная конкретная реализация:
\begin{description}[leftmargin=1.4em,itemsep=2pt]
\item[P1] Связная четырехмерная евклидова арена
      $\mathcal{M}^4$, $ds^2_E=\delta_{AB}dX^AdX^B$, без
      выделенной временной координаты.
\item[P2] $O(4)$-инвариантное фундаментальное действие,
      $S_E[\Phi]=S_E[\Lambda\cdot\Phi]$ при $\Lambda\in O(4)$.
\item[P3] Реальное скалярное поле $\Phi(X)$ с минимальным
      локальным микро-действием
      $S_E^{\rm bare}=\int d^4X\,[\tfrac{1}{2}(\partial_A\Phi)^2
        + V(\Phi)]$,
      $V=-\tfrac{\mu^2}{2}\Phi^2+\tfrac{\lambda}{4}\Phi^4$.
      Двухъямный потенциал здесь — это структура типа
      $\mathbb{Z}_2$, а не ``мексиканская шляпа''.
\item[P4] Физически релевантный вакуум принадлежит упорядоченной
      ветви с $\langle\partial_A\Phi\rangle=u_0\,\delta_{Aw}$,
      $u_0>0$, что выделяет предпочтительное направление $w$ и
      нарушает симметрию $O(4)\to O(3)$.
\item[P5] Поточечные внутренние координаты упорядоченной ветви
      (фаза $\theta(X)$, ориентация $n_A(X)$) являются
      описательными избыточностями; физические наблюдаемые зависят
      только от инвариантных градиентных и голономических данных.
      Переход от P5 к маршруту локальной калибровочной структуры,
      ассоциируемой с $SU(2)\times U(1)$, трактуется как
      производный структурный маршрут на гладкой когерентной
      ветви, а не как независимый постулат.
\item[P6] Множественность допустимых конфигураций и допустимость
      неравномерного упорядочения (домены, переходы, дефекты).
\end{description}

\textbf{Производное правило ветви (BR1).}
На гладком когерентном секторе упорядоченной ветви эффективный
параметр порядка имеет вид
$\Phi_{\rm eff}=\rho\,e^{i\theta}$, $\rho>0$.
Для замкнутых контуров, не пересекающих дефектные ядра,
\begin{equation}
  \label{eq:winding}
  \oint_{\mathcal C}\partial_A\theta\,dX^A = 2\pi n,
  \qquad n\in\mathbb{Z},
\end{equation}
а элементарная когерентная петля вводит доквантовую
топологическую шкалу действия
\begin{equation}
  \label{eq:loop_action}
  S_{\rm loop,min}=2\pi S_0,
\end{equation}
где $S_0$ является ECT-предшественником $\hbar$
(согласование: \statB).
BR1 является \emph{производным}: оно следует после ограничения на
гладкое однозначное коллективное фазовое поле на нарушенной фазе.

\textbf{P7 не входит в минимальную архитектуру.}
Постулат~P7 (комплексный внутренний модуль ранга~3 с эрмитовой
формой и фиксированной комплексной формой объема, дающий точную
локальную симметрию $SU(3)$) представляет собой отдельный
\emph{обратно сконструированный} слой, вводимый в части~IV
препринта. Он нужен потому, что $\mathfrak{su}(3)\not\subset
\mathfrak{so}(4)$: цветовой сектор не может жить внутри
одного только $O(4)$-упорядоченного конденсата. Следовательно,
P7 является добавкой к ядру P1--P6, а не частью самого ядра.

\textbf{Два энергетических масштаба и отдельная длина.}
Энергетическую размерность имеют $\phi_0\sim\bar M_{\rm Pl}$ и
$v_2\approx246$~GeV. Градиентная амплитуда
$u_0=|\langle\partial_A\Phi\rangle|$ имеет размерность GeV$^2$;
условная размерностная оценка имеет вид $u_0\sim\phi_0m_\sigma$.
Отдельно, внутри принятого Level-C галактического функционала,
$L_{\rm gal}=r_*=\sqrt{G_NM_{\rm bar}/g_\dagger}$ является длиной,
а не третьей энергетической амплитудой. ECT-мост к $L_{\rm gal}$ Open.
Логарифм $\mathcal I_{\rm EW}=\ln(\phi_0/v_2)\approx36.8$
организует Open-программу генерации электрослабого масштаба.

% =====================================================================
\section{Экономия входов и структурные переформулировки}
\label{sec:input_economy}
% =====================================================================

ECT расходует свои предположения на небольшой микроскопический
слой (\S\ref{sec:architecture}) и предлагает структурные маршруты
к нескольким секторам, которые в фундаментальной физике обычно
трактуются как независимые входы. Два компактных обзора ниже
заменяют длинную сводную таблицу: читатель, которому нужны детали,
найдет их в соответствующих разделах.

\begin{mechbox}
\textbf{Аудит входов.}
ECT не утверждает, что меньшее число предположений само по себе
доказывает теорию; задача в том, чтобы явно показать, где именно
предположения тратятся, а где заявляется реконструкция.
\begin{itemize}[leftmargin=1.2em,itemsep=1pt]
\item \textbf{Минимальный вход упорядоченной среды:}
      $\Phi$-среда на $\mathcal{M}^4$, вариационный принцип DP,
      рабочие постулаты P1--P6, производное правило ветви BR1.
\item \textbf{Восстановлено или структурно реконструировано:}
      лоренцев причинный конус; эффективная метрика и
      ориентационный сектор; когерентная квантовая ветвь;
      абелева решетка зарядов компактной группы $U(1)$; прямой канал
      космологической постоянной развязан.
\item \textbf{Согласовано / индуцировано к стандартной
      низкоэнергетической физике:}
      $c\equiv c_*$, $\hbar\equiv S_0$,
      $G_N=c_*^4/(8\pi M_G^2)$, энергетические масштабы
      $\phi_0,\,v_2$ и отдельная размерностно-квадратичная амплитуда
      $u_0$. Галактический $L_{\rm gal}$ является Level-C длиной fit-замыкания.
\item \textbf{Слои замыкания и диагностики:}
      сохраняемое космологическое $\varepsilon$-окно и
      поздневременное инфракрасное масштабирование. Галактический
      закон среды/истории $\Xi$ — Open-гипотеза для будущего теста,
      а не реализованное замыкание.
\item \textbf{Постулируется отдельно:}
      P7 (комплексный внутренний модуль ранга~3) для точной
      локальной симметрии $SU(3)$.
\item \textbf{Открыто:}
      динамическое происхождение самой упорядоченной ветви
      (градиентная конденсация, OP-grad);
      полная теорема Борна из первых принципов и проблема
      единственного исхода и правила обновления при измерении;
      независимый вывод предпосылок редуцированного описания
      C1--C3 (\S\ref{sec:quantum});
      Bell/CHSH и предел Цирельсона;
      полная низкоэнергетическая QFT-архитектура (локальность,
      кластерное разложение, аналитичность, симметрия кроссинга,
      S-матрица);
      наблюдаемое значение $\rho_\Lambda$;
      гиперзаряд и полный спектр заряда/ароматов/поколений
      Стандартной модели;
      strong-CP проблема в CP-нечетной топологической связи;
      завершенное УФ-завершение.
\end{itemize}
\end{mechbox}

\textbf{Структурные переформулировки, а не окончательные решения.}
ECT не подает перечисленные ниже пункты как решенные на одном и
том же уровне. Важно увидеть, где та же архитектура упорядоченной
среды дает структурную переформулировку, а где сохраняются
открытые зависимости:

\begin{itemize}[leftmargin=1.4em,itemsep=2pt]
\item \textbf{Проблема времени.}
      Время возникает как направление упорядоченной ветви и как
      эффективный параметр на нарушенной фазе; полная
      квантово-гравитационная проблема времени автоматически не
      решается.
\item \textbf{Квантово-классическая граница.}
      Когерентный против декогерированного/ориентационного
      сектора как структурный критерий; строгая проблема
      измерения / единственного исхода остается \statO.
\item \textbf{Правило Борна.}
      Условный маршрут \statB{} через RP/OS-маршрут и аргумент
      Глисона при
      предпосылках C1--C3 (\S\ref{sec:quantum});
      вывод C1--C3 из голых P1--P6 и проблема
      единственного исхода и правила обновления остаются \statO.
\item \textbf{Катастрофа космологической постоянной.}
      Прямой канал фоновой вакуумной энергии развязывается на уровне
      \statA{} благодаря фиксированному детерминанту кинетики
      упорядоченной ветви; наблюдаемое значение $\rho_\Lambda$
      остается \statC/\statO.
\item \textbf{Галактический темный сектор.}
      Принятая интерполяция дает Level-C fit для RAR/кривых
      вращения; наклон BTFR~4 — условная алгебра B/C внутри этого
      функционала.  Скалярно-галактический мост, $\Xi$ и
      UDG/cluster-интерфейс остаются \statO.
\item \textbf{Информация черных дыр.}
      Критическая оболочка / редуцированное состояние как
      отображение на язык островов и квантовых экстремальных
      поверхностей;
      полный расчет кривой Пейджа внутри ECT остается \statO.
\item \textbf{Калибровочный / цветовой сектор.}
      Структурные маршруты абелевой решетки компактной группы
      $U(1)$ и
      хирального $SU(2)$ из нарушенной фазы; точный локальный
      $SU(3)$ — через отдельный постулат P7; полный спектр
      зарядов/ароматов/поколений Стандартной модели, конфайнмент и
      strong-CP проблема остаются \statO.
\end{itemize}

Смысл этого раздела не в том, чтобы доказать ECT за счет
лаконичной подачи и не в том, чтобы предъявить список уже
разрешенных парадоксов. Его задача — явно показать, где ECT
тратит предположения и где именно она заявляет реконструкцию;
статусные метки остаются управляющим стандартом.

% =====================================================================
\section{Дисциплина статусов утверждений}
\label{sec:status_discipline}
% =====================================================================

Каждый количественный или качественный результат в ECT помечен
одним из четырех статусных уровней. Эта таксономия используется
во всем препринте и ниже применяется единообразно.

\begin{description}[leftmargin=1.4em,itemsep=2pt]
\item[\statA] Строгий структурный вывод или теорема в рамках
      явно сформулированных предпосылок. Пример: эмерджентное
      унимодулярное развязывание прямого канала фоновой
      вакуумной энергии;
      возникновение лоренцева конуса внутри нарушенной фазы при
      данных условиях анизотропной жесткости.
\item[\statB] Условный или полученный на уровне замыкания вывод;
      структурный
      маршрут существует, но для получения числа требуется хотя бы
      один шаг согласования, калибровки или перехода к
      эффективному функционалу.
      Пример: Level-B согласование $S_0\leftrightarrow\hbar$; наклон
      BTFR как условная B/C-алгебра внутри принятого Level-C
      функционала G; B/C comparison ansatz
      $g^\dagger_0\sim cH_0/(2\pi)$.
\item[\statC] Феноменологический или диагностический слой
      согласованности. Пример: сохраняемое пятипробное эффективное
      окно $\varepsilon\in[0.0296,0.0376]$ при $1\sigma$.
      Это эмпирическая совместимость, а не предсказание из первых
      принципов.
\item[\statO] Явно неразрешенный пункт. Пример: строгая теорема
      Борна из первых принципов и проблема единственного исхода;
      наблюдаемое значение $\rho_\Lambda$; полная структура
      зарядов, ароматов и поколений Стандартной модели;
      strong-CP проблема в CP-нечетном топологическом секторе;
      закон среды $\Xi$.
\end{description}

Дисциплина статусов намеренно асимметрична: ослабить результат с
\statA{} до \statB{} допустимо, если какой-то шаг вывода окажется
неполным; усиливать в обратную сторону нельзя без явного строгого
доказательства.

% =====================================================================
\section{Логика вывода и карта механизмов}
\label{sec:logic}
% =====================================================================

Архитектура из \S\ref{sec:architecture}, аудит входов из
\S\ref{sec:input_economy} и дисциплина статусов из
\S\ref{sec:status_discipline} вместе задают единый поток
механизмов, которому и посвящена остальная часть этого обзора.
На рисунке~\ref{fig:logic} он дан в компактном виде.

\begin{figure}[h]
\centering
\begin{tikzpicture}[
  node distance=10mm and 8mm,
  every node/.style={font=\footnotesize},
  box/.style={draw, rounded corners=2pt, align=center,
              inner sep=4pt, minimum width=3.9cm,
              minimum height=1.0cm, fill=gray!4},
  side/.style={draw, rounded corners=2pt, align=center,
              inner sep=4pt, minimum width=5.0cm,
              minimum height=0.9cm, fill=gray!14},
  arr/.style={->, thick, >=stealth},
  outarr/.style={->, thick, >=stealth, dashed}
]

\node[box] (medium)
  {$\Phi$-среда на $\mathcal{M}^4$\\
   (нет времени, метрики и квантов)};

\node[box, below=of medium] (ordered)
  {$O(4)\!\to\!O(3)$ градиентное упорядочение\\
   \statA{} (структурный механизм)};

\node[box, below=14mm of ordered] (geom)
  {Скалярное замыкание\\
   ориентация Open};

\node[box, left=of geom] (coh)
  {Когерентная ветвь\\
   $S_{\rm loop,min}=2\pi S_0$\\
   $S_0\leftrightarrow\hbar$ (Level-B matching), OS/RP-мост\\
   сектор Шредингера};

\node[box, right=of geom] (ir)
  {Галактический fit \statC\\
   BTFR \statB/\statC; $\Xi$ \statO};

\draw[arr] (medium) -- (ordered);
\draw[arr] (ordered.south) -- ++(0,-3mm) -| (coh.north);
\draw[arr] (ordered) -- (geom);
\draw[arr] (ordered.south) -- ++(0,-3mm) -| (ir.north);

\node[side, below=16mm of geom] (p7)
  {Отдельный боковой слой: постулат P7\\
   комплексный внутренний модуль ранга~3};

\node[side, below=of p7] (su3)
  {Точное локальное $SU(3)$ цветовое завершение\\
   P7 постулируется; полная феноменология \statO};

\draw[outarr] (p7) -- (su3);

\end{tikzpicture}
\caption{Упрощенная логика вывода в ECT. Стрелки показывают поток
механизмов, а не утверждения одинаковой теоремной силы. Ветвь
$SU(3)$ вынесена отдельно, потому что требует P7 и не следует из
минимальной $O(4)$-упорядоченной среды; веса типа Борна относятся
к условному слою квантовой реконструкции, обсуждаемому в
\S\ref{sec:quantum}.}
\label{fig:logic}
\end{figure}

Тот же поток можно выразить словами как пять структурных путей от
одной среды к качественно различным явлениям:

\begin{enumerate}[leftmargin=1.4em,itemsep=2pt]
\item \textbf{Упорядоченная ветвь} (P4, анизотропный кинетический
      тензор с $\alpha>\beta$): выделение лоренцева причинного
      конуса и эффективной сигнатуры $(-,+,+,+)$.
      \S\ref{sec:cone}.
\item \textbf{Гладкий когерентный сектор}
      (BR1, гладкая однозначная $\theta$): целочисленная обмотка,
      $S_{\rm loop,min}=2\pi S_0$, отражательная положительность
      и мост Остервальдера--Шрадера, дающий
      \emph{частичный маршрут реконструкции гильбертова
      пространства} для когерентного квантового сектора.
      \S\ref{sec:quantum}.
\item \textbf{Скалярный метрический отклик и открытое
      ориентационное расширение}: показанное уравнение является
      скалярным.  Физический $T^{(n)}_{\mu\nu}$ требует ковариантного
      $S_n[g,n,\lambda_n]$; $\mathcal X^{\rm cand}_{\mu\nu}$ — лишь
      внешний прокси.  \S\ref{sec:gravity}.
\item \textbf{Заданная галактическая интерполяция}: Level-C
      fit для RAR/кривых вращения и условная алгебра BTFR.  Она не
      выводится из однопетлевого результата для $\psi_\chi$ или из
      скалярного трейс-уравнения.  Для космологии отдельно используется
      скалярное диагностическое замыкание.
      \S\S\ref{sec:cosmology}--\ref{sec:galactic}.
\item \textbf{Обратно сконструированный цветовой слой}
      (P7, комплексный внутренний модуль ранга~3): точное
      локальное $SU(3)$ ценой отдельного постулата, выходящего за
      пределы $O(4)$-упорядоченного конденсата.
      \S\ref{sec:gauge}.
\end{enumerate}

% =====================================================================
\section{Чем ECT отличается от близких подходов}
\label{sec:comparison}
% =====================================================================

ECT не является простым перекомбинированием уже существующих
программ. Две таблицы ниже размещают ECT относительно ближайших
устоявшихся или предлагаемых теоретических рамок; разбиение на
\emph{концептуальные основания}
(таблица~\ref{tab:ect_vs_nearby_A}) и
\emph{феноменологию и калибровочный сектор}
(таблица~\ref{tab:ect_vs_nearby_B}) введено только ради
читаемости.

\renewcommand{\arraystretch}{1.22}
\small
\footnotesize
\setlength{\tabcolsep}{3.5pt}
\begin{longtable}{@{}L{2.4cm} L{2.8cm} L{2.6cm} L{4.3cm} L{2.2cm}@{}}
\caption{Концептуальные основания.}
\label{tab:ect_vs_nearby_A}\\
\toprule
\textbf{Сектор} & \textbf{Стандартный подход} &
\textbf{Близкие альтернативы} & \textbf{Механизм ECT} &
\textbf{Статус}\\
\midrule
\endfirsthead
\multicolumn{5}{@{}l}{\textit{Таблица~\ref{tab:ect_vs_nearby_A},
продолжение.}}\\
\toprule
\textbf{Сектор} & \textbf{Стандартный подход} &
\textbf{Близкие альтернативы} & \textbf{Механизм ECT} &
\textbf{Статус}\\
\midrule
\endhead
\bottomrule
\endfoot

Сигнатура пространства-времени
& Лоренцева метрика предполагается изначально
& CDT~\cite{Ambjorn2005}, causal sets, analogue gravity;
  маршрут через clock-field~\cite{MukohyamaUzan2013}
& Упорядоченная евклидова ветвь выделяет эффективный лоренцев
  конус; формула и согласование приведены в \S\ref{sec:cone}.
& \statA{} существование; \statB{} согласование\\

Квантовый формализм
& Гильбертово пространство и правило Борна постулируются
& many-worlds, Bohm, QBism, OS-реконструкция
& Когерентная ветвь дает маршрут RP/OS к гильбертовой динамике;
  веса типа Борна остаются условными.
& \statB{} структурно; теорема Борна — \statO\\

Гравитация
& Метрическое поле фундаментально
& индуцированная гравитация Сахарова, термодинамический вывод
  Джейкобсона~\cite{Jacobson1995},
  Einstein--aether/vector-tensor аналоги~\cite{Jacobson2004},
  энтропийная гравитация Верлинде~\cite{Verlinde2011}, analogue
  gravity~\cite{Volovik2003,Padmanabhan2010}
& Скалярный анзац; ориентация Open.
& \statB{}; \statO{}\\

Космологическая постоянная
& Вакуумная энергия гравитирует; требуется сокращение
  порядка $\sim 10^{120}$
& унимодулярная гравитация, sequestering, аргументы no-go
  Вайнберга
& Прямой канал фоновой вакуумной энергии развязывается через фиксированный
  детерминант упорядоченной ветви; наблюдаемое значение остается
  диагностическим / открытым.
& \statA{} для прямого канала; наблюдаемое $\rho_\Lambda$
  — \statC/\statO\\

\end{longtable}
\normalsize

\textbf{Независимая близкая работа по возникновению лоренцевой
сигнатуры.}
Mukohyama и Uzan исследовали возникновение лоренцевой динамики из
четырехмерной римановой микроскопической структуры через фон
производной поля-часов~\cite{MukohyamaUzan2013}. Это близкая
предшествующая работа по emergent Lorentz signature, но она не
является основанием ECT: настоящая конструкция использует иную
архитектуру скалярной среды, упорядоченный градиентный конденсат,
разделение на когерентную и декогерированную ветви и иную
гравитационную и космологическую программу.

\renewcommand{\arraystretch}{1.22}
\small
\footnotesize
\setlength{\tabcolsep}{3.5pt}
\begin{longtable}{@{}L{2.4cm} L{2.8cm} L{2.6cm} L{4.3cm} L{2.2cm}@{}}
\caption{Феноменология и калибровочный сектор.}
\label{tab:ect_vs_nearby_B}\\
\toprule
\textbf{Сектор} & \textbf{Стандартный подход} &
\textbf{Близкие альтернативы} & \textbf{Механизм ECT} &
\textbf{Статус}\\
\midrule
\endfirsthead
\multicolumn{5}{@{}l}{\textit{Таблица~\ref{tab:ect_vs_nearby_B},
продолжение.}}\\
\toprule
\textbf{Сектор} & \textbf{Стандартный подход} &
\textbf{Близкие альтернативы} & \textbf{Механизм ECT} &
\textbf{Статус}\\
\midrule
\endhead
\bottomrule
\endfoot

Галактический темный сектор
& Частицы темной материи, NFW-гало
& MOND/AQUAL~\cite{Milgrom1983,Famaey2012}, emergent dark gravity
  Верлинде
& Заданный интерполяционный fit; наклон BTFR~4 является
  условной алгеброй глубокой ветви (\S\ref{sec:galactic}).
& fit \statC; наклон \statB/\statC; $\Xi$ \statO\\

Напряжения Хаббла / $S_8$
& Расширения набора параметров $\Lambda$CDM, EDE, модифицированная
  гравитация
& EDE, interacting DE, $f(R)$~\cite{Sotiriou2010},
  quintessence~\cite{Carroll2001}
& Один поздневременной параметр деформации $\varepsilon$;
  сохраняемое пятипробное окно согласованности на уровне
  $1\sigma$ (\S\ref{sec:cosmology}).
& \statC{} как слой совместимости; вывод из первых принципов
  — \statO\\

Калибровочный сектор / $SU(3)$
& Группа просто задается Стандартной моделью
& GUT, программы внутренней геометрии
& Выведенная абелева решетка зарядов компактной группы $U(1)$;
  хиральное $SU(2)$ из ориентированной нарушенной фазы;
  точный локальный $SU(3)$ через отдельный постулат P7.
& $U(1)$ — \statA/\statB; $SU(2)$ — \statB;
  $SU(3)$ через P7; полный спектр СМ — \statO\\

\end{longtable}
\normalsize

Еще три различия стоит сформулировать явно:
ECT \emph{не} предполагает существование лоренцева
пространства-времени как фундаментального ингредиента;
заданный галактический fit не содержит члена, описывающего гало частиц, но
не исключает частицы скрытого сектора; поздневременное замыкание также не
вводит темную энергию как внешнюю жидкость;
и ECT \emph{не} решает strong-CP проблему.

% =====================================================================
\section{Эмерджентное пространство-время и универсальный конус}
\label{sec:cone}
% =====================================================================

В нарушенной фазе линеаризованные возбуждения вокруг
упорядоченного фона распространяются с эффективным кинетическим
тензором, который становится анизотропным в направлении $w$,
выделенном величиной $\langle\partial_A\Phi\rangle$.
Для канонических коэффициентов анизотропной жесткости
$(\alpha,\beta)$ продольно-поперечное расщепление допускает окно
лоренцевой сигнатуры тогда и только тогда, когда $\alpha>\beta$;
в этом случае эффективный низкоэнергетический конус имеет одну
универсальную скорость распространения:
\begin{equation}
  \label{eq:cstar}
  c_* = \sqrt{\beta/(\alpha-\beta)},
  \qquad
  \text{лоренцево окно: } \alpha>\beta.
\end{equation}

\begin{mechbox}
\textbf{Механизм в ECT.}
SSB $O(4)\!\to\!O(3)$ градиентного конденсата выделяет
предпочтительное направление; кинематическое ограничение
$n^An_A=1$ на ориентационную переменную в сочетании с
$\alpha>\beta$ заставляет продольную моду диспергировать на одном
эмерджентном световом конусе. Согласование $c_*$ с наблюдаемой
вакуумной скоростью $c$ — отдельный калибровочный шаг.\\
\textbf{Что этот результат устанавливает.} Явный маршрут от
евклидова упорядочения к лоренцевой сигнатуре, эффективному
временному направлению и универсальному низкоэнергетическому
причинному конусу внутри нарушенной фазы.\\
\textbf{Статус.} \statA{} для возникновения конуса внутри
нарушенной фазы; \statB{} для отождествления $c_*$ с наблюдаемым
$c$.\\
\textbf{Где это в полном препринте.} Часть~I, разделы по
упорядоченной ветви и классификации метрик; теорема об
универсальности конуса.\\
\textbf{Фальсификатор.} Неподавленная секторно-зависимая
низкоэнергетическая скорость распространения
(например, различие фотонной и гравитонной скоростей за пределами
ограничения GW170817) фальсифицировала бы возникновение
универсального конуса.
\end{mechbox}

% =====================================================================
\section{Квантовый сектор: что восстановлено и что остается открытым}
\label{sec:quantum}
% =====================================================================

ECT реконструирует квантовый сектор на когерентной ветви,
а не постулирует гильбертово пространство на фундаментальном
уровне.
Гладкий когерентный сектор упорядоченной ветви допускает
частичную реконструкцию стандартного квантового формализма через
евклидов интеграл по траекториям, где отражательная
положительность обеспечивает мост Остервальдера--Шрадера к
гильбертову
пространству и унитарной эволюции в стационарном когерентном
секторе. На уровне обзора эволюция в когерентном секторе
записывается схематически как
\begin{equation}
  \label{eq:qs}
  iS_0\,\partial_t\Psi
  = \hat H_{\rm coh}\Psi
  + \Delta_{\rm dec}\Psi
  + \Delta_{\nabla}\Psi,
\end{equation}
где $\Delta_{\rm dec}$ обозначает поправки редуцированного
состояния / декогеренции из функционала влияния, а
$\Delta_{\nabla}$ собирает старшие градиентные поправки к
кинетическому оператору; явная форма этих операторов приведена в
соответствующей части полного препринта и не нужна на уровне
обзора. Для стационарных когерентных собственных мод,
$\hat H_{\rm coh}\Psi=E\Psi$, $\Psi\sim e^{-i\omega t}$,
получается $E=S_0\omega$ с учетом контролируемых поправок.
Стандартное соотношение Планка--Эйнштейна $E=\hbar\omega$ затем
следует из \statB{}-согласования $S_0=\hbar$.

\textbf{Веса типа Борна как условная реконструкция.}
ECT не вставляет веса типа Борна в конструкцию аксиоматически.
Вместо этого, как только редуцированное
когерентно/декогерированное описание достигает трех конкретных
предпосылок — C1: декогерированные альтернативы допускают
эффективную гильбертову / проекторную структуру; C2: эти
альтернативы приблизительно ортогональны; C3: вероятности
складываются по взаимоисключающим альтернативам, — гильбертова
RP/OS-конструкция вместе с аргументом уникальности типа Глисона
выделяет веса Борна. Такая условная реконструкция имеет статус
\statB. Открытым остается более ранний шаг: вывод C1--C3
непосредственно из голых постулатов ECT, получение теоремы Борна
из первых принципов без этих предпосылок редуцированного
описания и решение проблемы единственного исхода / правила
обновления при измерении. Common-medium nonfactorisation является кандидатным маршрутом
\statB{}/\statO{}; физический state/channel, Bell/no-signalling
корреляторы и предел Цирельсона остаются \statO.

Связь спин--статистика установлена на уровне \statA{} для
свободного когерентного ядра через евклидов маршрут
Остервальдера--Шрадера (аргумент класса
L\"uscher--Osterwalder, сверенный с исходными источниками);
завершение с взаимодействиями остается \statO.
Свободное антисимметричное ядро даёт точное математическое Pauli/CAR
утверждение; применение к физической взаимодействующей ECT-материи
остаётся \statB{}/\statO{} до spinor/kernel completion.
Условный маршрут к CPT идентифицирован через евклидову структуру
Остервальдера--Шрадера; оператор выделенного направления
появляется лишь как CPT-нечетный оператор, чувствительный к
выбору ветви,
совместимый с точным симметрийным каркасом.
PES-R даёт статусно разделённый инвентарь: action stationarity; pairwise
record exponent $\Phi_{ab}(T)$; declared translation diagnostic
$\Gamma_{\rm tr}[q;\eta]$; sector width $\gamma_a$ (или
$S_0\gamma_a$ в action units); и protocol entropy production
$\Sigma=\ln(P_F/P_R)$ только для именованного forward/reverse
протокола. Он не является универсальной теоремой отбора
persistent-сектора; additive mean-EOM functional, intersection theorem
и PES-to-Born bridge не выведены.

\begin{mechbox}
\textbf{Механизм в ECT.}
Отражательная положительность и OS-реконструкция на когерентной
ветви вместе с топологической доквантовой шкалой
$S_{\rm loop,min}=2\pi S_0$.\\
\textbf{Что этот результат устанавливает.} Маршрут от
когерентной ветви к гильбертовой динамике, сектору Шредингера и
топологическому предшественнику $\hbar$; веса типа Борна
возникают после введения предпосылок C1--C3.\\
\textbf{Статус.} \statB{} для OS-моста, PES и согласования
$S_0=\hbar$; \statB{} условно для весов типа Борна при принятии
C1--C3; \statO{} для теоремы Борна из голых P1--P6, для самих
C1--C3 и для проблемы единственного исхода / обновления при
измерении; \statA{} для связи спин--статистика на свободном
когерентном ядре (завершение с взаимодействиями — \statO).\\
\textbf{Где это в полном препринте.} Часть~III, разделы о
когерентной динамике, PES, гибридной согласованности и модели
декогеренции через функционал влияния.\\
\textbf{Дискриминант.} Надёжное нарушение явно согласованной
когерентной динамики ограничило бы это замыкание. Нынешние
gravity-decoherence и mediator-loophole bounds проверяют внешние
comparator classes, а не абсолютную скорость ECT: ECT record vertex,
state и noise kernel остаются Open/Not Identifiable.
\end{mechbox}

% =====================================================================
\paragraph{DP-02: отрицательный спектральный результат.}
Проверенный минимальный скалярный вакуумный канал имеет
$J_\phi\propto\omega^3$; equal-mass Gaussian composite vacuum —
$J_{\mathcal O}\propto\omega^9$ с finite-time огибающей $T^{-8}$; тепловой
режим — $J_{\mathcal O}^{\rm th}\propto T_b^5/k$. Ни один не даёт
универсальную инфракрасную форму Диоси--Пенроуза $1/k^2$. Альтернативные
ECT state/pole/vertex конструкции остаются Open.

\section{Гравитация и катастрофа космологической постоянной}
\label{sec:gravity}
% =====================================================================

Контролируемое метрическое уравнение показанного
скалярного анзаца в системе Йордана в неразделенной конвенции имеет вид
\begin{equation}
  \label{eq:gee}
  F\,G_{\mu\nu}
  = \frac{8\pi G_N}{c^4}\,T_{\mu\nu}
    + \mathcal{T}^{(q)}_{\mu\nu},
\end{equation}
где $\mathcal{T}^{(q)}_{\mu\nu}$ собирает производные и
потенциальные члены скаляра.  Множитель $G_N/F$ возникает лишь
после деления \emph{всего} уравнения на $F$.  Показанное действие не
содержит независимого ориентационного поля и потому не выводит
$T^{(n)}_{\mu\nu}$; $\mathcal X^{\rm cand}_{\mu\nu}$ является только
внешним прокси.  Ковариантное $S_n[g,n,\lambda_n]$ и его метрическая
вариация остаются \statO.  Деление уравнения на $F$ не создает
ориентационного члена.  Скалярный аудит использует
$a=f_{,q}/f$, $\mathcal A=K+3f_{,q}^2/(2f)$ и
$\mathcal E=V_{,q}-2aV-(a/2)\rho_{\rm m}$; при устойчивом однородном
корне $d\ln\bar f/d\rho_{\rm m}=a^2/(2\mathcal E_{,q})\ge0$ и
$m_{\rm loc}^2=\mathcal E_{,q}/\mathcal A$ являются условными локальными
результатами, а не экранированием конечного тела.
FRW/Hubble/JWST-диагностика использует скалярный фон.  Планковское согласование
$G_N=c_*^4/(8\pi M_G^2)$ сохраняется как отдельное соотношение.

M (микроскопический director-оператор), S (этот Level-B скалярный анзац) и G (принятый Level-C галактический функционал) различны. M$\neq$S по операторному составу, а S$\to$G остаётся Open. Историческое LOF-приложение является смешанным аудитом ветвей: результат для тяжёлого радиального прокси не выбирает другую координату или G, а matter-to-Lorentz-order map остаётся Open.

\textbf{Катастрофа космологической постоянной (прямой канал).}
В стандартной общей релятивистской рамке базовый вакуумный вклад
четвертичного двухъямного потенциала
$V(\phi_0)=-\mu^4/(4\lambda)$ гравитировал бы как
космологическая постоянная порядка $\lambda M_{\rm Pl}^4/4$,
что превышает $\rho_\Lambda^{\rm obs}
\sim(2.3\times 10^{-3}\,{\rm eV})^4$ примерно на $10^{120}$.
В ECT эта катастрофа устраняется на уровне \statA{} для ведущего
сектора упорядоченной ветви благодаря теореме об
\emph{эмерджентном унимодулярном развязывании}: детерминант
кинетического тензора упорядоченной ветви фиксируется
кинематическим ограничением $n^An_A=1$:
\begin{equation}
  \sqrt{-\det K^{AB}} = \beta^2/c_*,
\end{equation}
и, как следствие, голый вклад $V(\phi_0)$ не источает локальную
эмерджентную гравитационную динамику. При дополнительном
однострочном спектральном предположении о полном квадратичном
операторе флуктуаций (условие $H_\Lambda$) нулевая по производным
однопетлевая вакуумная добавка поглощается той же фиксированной
структурой кинетического детерминанта. Поэтому прямое сокращение
ультрафиолет-инфракрасного масштаба порядка $10^{120}$,
характерное для стандартной схемы, структурно отсутствует в
ведущем секторе упорядоченной ветви.
В текущем препринте зафиксированы и два дальнейших усиления:
утверждение о развязке реализовано на явном точном FRW-фоне
упорядоченной ветви (задачи реализации в программе
космологической постоянной завершены), а также установлена
голографическая энтропийная оценка area-law как структурное
свойство упорядоченной среды
(\statA{} для масштабирования, \statB{} для коэффициента).

\textbf{Отдельная оставшаяся задача.}
Наблюдаемое численное значение $\rho_\Lambda$ остается отдельной
задачей уровня \statC; естественное инфракрасное масштабирование
$\rho_\Lambda\sim c_\Lambda M_{\rm Pl}^2 H_0^2$ при
$c_\Lambda=\mathcal{O}(1)$ согласуется с наблюдениями по порядку
величины, но сам коэффициент пока не выведен.

\begin{mechbox}
\textbf{Механизм в ECT.}
Скалярный метрический отклик внутри объявленного анзаца Йордана;
фиксированный кинетический детерминант через $n^An_A=1$ относится
к отдельному результату о развязке прямого канала.\\
\textbf{Что этот результат устанавливает.} Согласованную
нормировку скалярного метрического уравнения и развязку прямого
канала фоновой вакуумной энергии.  Ковариантный ориентационный
стресс показанным действием не установлен.\\
\textbf{Статус.} точная алгебра нормировки внутри объявленного
Level-B скалярного анзаца; отдельно \statA{} для теоремы о
развязке прямого канала; \statO{} для ориентационного
расширения; \statC{} для наблюдаемого значения $\rho_\Lambda$;
\statO{} для первопринципного значения $c_\Lambda$.\\
\textbf{Где это в полном препринте.} Часть~II, разделы по
обобщенным уравнениям Эйнштейна, проблеме космологической
постоянной и теореме о развязке.\\
\textbf{Фальсификатор.} Надежная демонстрация того, что постоянные
сдвиги вакуумной энергии источают локальные эмерджентные
гравитационные уравнения со стандартной силой GR,
фальсифицировала бы утверждение о развязке прямого канала.
Независимо от этого статистически устойчивый переход в область
$w<-1$ оказал бы сильное давление на текущую схему замыкания
темной энергии на упорядоченной ветви.
\end{mechbox}

% =====================================================================
\section{Космология и текущие наблюдательные напряжения}
\label{sec:cosmology}
% =====================================================================

Макроскопическое замыкание упорядоченной ветви допускает
однородную эффективную поздневременную деформацию,
\begin{equation}
  \label{eq:eps_def}
  G_{\rm eff}(z)/G_N = (1+z)^{2\varepsilon},
\end{equation}
где $\varepsilon$ — единый диагностический параметр
\emph{сохраняемого слоя анализа по пяти пробам}
(Hubble плюс sound horizon $r_s$, извлечение JWST при больших $z$,
cosmic chronometers, красносмещенные искажения $f\sigma_8$ и
integrated Sachs--Wolfe cross-correlation). Общее эффективное
ограничение по этим пяти сохраненным пробам дает
\begin{equation}
  \varepsilon\in[+0.0296,\,+0.0376] \quad (1\sigma), \quad
  [+0.021,\,+0.042] \quad (2\sigma).
\end{equation}
Это \emph{окно согласованности} между наблюдательно различными
каналами с разными ядрами и систематиками.
Его следует читать как косвенный многоканальный
\emph{тест жизнеспособности} текущей эффективной параметризации,
а не как уже завершенную глобальную статистическую комбинацию или
замкнутый вывод из первых принципов.

Связанные с лоренцевой ветвью индикаторы возраста составляют
$\sim 13.5$~Gyr при отображении с базовым $H_0$ и
$\sim 12.5$~Gyr при отображении с выводимым $H_0$; сектор ранних
галактик JWST дает связывающую нижнюю границу сохраняемой полосы
в текущем конвейере извлечения. Отдельно перечислены исключенные
пробы: BAO и линзирование $A_{\rm lens}$ как методологически
ограниченные в рамках однородного анзаца для $\varepsilon$;
$S_8$~(KiDS-1000) и $\sigma_8$ по скоплениям как находящиеся
вне текущего $\varepsilon$-сектора при интерпретации через
категориальную ошибку; канал SN~Ia как не независимый относительно
сохраненных якорей.

В рамках того же фона упорядоченной ветви поздневременной сектор
темной энергии представляет собой \emph{thawing}-замыкание.
Точное фоновое соотношение имеет вид
$w_0=(r_0-1)/(r_0+1)$, где $r_0$ — современное отношение
кинетической энергии к энергии конденсата
($r_0\approx 0.095$ в калиброванной реперной точке),
с линеаризованной формой
$w(z)\simeq -1+2f_0\,e^{-\kappa z}$,
$f_0\approx 0.087$.
Часто цитируемое значение $w_0\approx -0.83$ — это
реперная \emph{thawing}-точка, откалиброванная по DR1
(\statB): полезная реперная точка, а не вывод прямо из
микроскопического слоя.
Внутри этого замыкания выполняется неравенство $w(z)\ge -1$ при
всех красных смещениях.

\begin{mechbox}
\textbf{Механизм в ECT.}
Единый параметр деформации упорядоченной ветви $\varepsilon$,
отслеживающий тот же конденсатный фон, который контролирует и
сектор темной энергии, и зависящий от эпохи базовый масштаб
ускорения
$g^\dagger_{\rm bg}(z)=c_*H(z)/(2\pi)$.\\
\textbf{Что этот результат устанавливает.} Сохраняемое
поздневременное окно согласованности, в котором один эффективный
параметр деформации организует сразу несколько наблюдательных
каналов.\\
\textbf{Статус.} \statC{} как слой совместимости; вывод
$\varepsilon(z)$ из функций замыкания из первых принципов остается
\statO.\\
\textbf{Где это в полном препринте.} Часть~II, раздел по
сохраняемому анализу пяти проб и его методологическое приложение.\\
\textbf{Фальсификатор.} Будущий совместный анализ тех же
сохраненных проб, который устойчиво уводит $\varepsilon$ ниже
нуля (в особенности если извлечение JWST при больших $z$
стабилизируется на связывающей нижней границе
$\varepsilon\le 0$), оказал бы серьезное давление на текущий
сектор.
\end{mechbox}

% =====================================================================
\section{Галактическая феноменология}
\label{sec:galactic}
% =====================================================================

В глубоком режиме малых ускорений для принятого галактического замыкания
асимптотическое галактическое уравнение поля принимает вид
$g\simeq\sqrt{g_N\,g_\dagger}$, а скорость на асимптотически
плоской кривой вращения удовлетворяет
\begin{equation}
  \label{eq:btfr}
  v_\infty^4 = G_N\,M_b\,g_\dagger.
\end{equation}
Асимптотический наклон BTFR, равный $4$, является алгебраическим
следствием только внутри заданной интерполяции глубокой ветви; само
замыкание не выведено из исправленного скалярного действия.  Координата
$\phi_u=\beta_\phi^{-1}\ln(u/u_\infty)$ отлична от
$\psi_\chi=\ln(\chi/\chi_{\rm vac})$, а предварительный результат
$K_\chi\propto e^{-\psi_\chi/2}$ не выводит AQUAL-ветвь.
Практический fit возвращает отдельные значения $g^\dagger_{\rm eff}$ для галактик (Level~C).  Референс $g^\dagger_0\approx1.04\times10^{-10}$\,м/с$^2$ является согласованным B/C-бенчмарком, сравниваемым с эмпирическим $a_0\approx1.2\times10^{-10}$\,м/с$^2$, а не универсальным результатом fit.  Возможная регрессия $g^\dagger_{\rm eff}/g^\dagger_0=\Xi(\rho_{\rm env},\bar\phi_{\rm env},\ldots)$ — только Open-гипотеза для будущего теста.  Медианы выборки равны $\chi_r^2\approx2.7$ при фиксированном $M/L$ и $\approx1.4$ при свободном $M/L$, против quoted MOND comparator $\approx3.1$ на сопоставимом числе параметров.  Условная длина $L_{\rm gal}=r_*=\sqrt{G_NM_b/g^\dagger}$ принадлежит G и не является энергетическим масштабом конденсата. Мост S$\to$G, фундаментальный состав частиц и fraction budget остаются Open.

Исправленный UDG inverse diagnostic использует коэффициент Wolf~3 при
$R_{1/2}=4R_e/3$ и неограниченную алгебраическую величину
$\Xi_{\rm alg}=\mathcal R_{\rm dyn}(\mathcal R_{\rm dyn}-1)g_N/g^\dagger_0$.
Она может интерпретироваться как физический environment factor только при
$\Xi_{\rm alg}\ge0$. Центральные значения: $-0.00369$ для DF4 (нет
допустимого решения $\Xi\ge0$), $0.00516$ для FCC~224, $0.892$ для
NGC~5846-UDG1 и $4.981$ для Dragonfly~44; dispersion bracket DF2 даёт
$0.031$--$0.423$. Прежняя логарифмическая 95\%-полоса nuisance-параметров
снята; полный Jeans posterior остаётся Open.

Кривая force dimensionality относится к отдельному иллюстративному
point-source закону $\mu_{\rm dim}=x/\sqrt{1+x^2}$, а не к практическому G;
проверка протяжённых галактик требует измеренного $p_N(r)$ и не выполнена.
Кластерный скрипт устанавливает лишь точное сохранение порядка/argmax для
вручную заданных карт. Это не морфологическое или линзинговое свидетельство
и не наблюдаемая amplitude inference; физический response и likelihood Open.

\begin{mechbox}
\textbf{Механизм в ECT.}
Принятое низкоускорительное интерполяционное замыкание;
наклон BTFR как алгебраическое следствие его глубокой ветви;
возможная будущая регрессия per-galaxy $g^\dagger_{\rm eff}$ через заранее заданную Open-функцию $\Xi$; текущие fit её не используют.\\
\textbf{Что этот результат устанавливает.} Феноменологический
fit-бенчмарк и условный алгебраический наклон BTFR внутри принятого
замыкания; физический мост $u\leftrightarrow\chi$ и скалярное
экранирование этим не устанавливаются.\\
\textbf{Статус.} Наклон BTFR — \statB{}/\statC{} условно на замыкании; per-galaxy $g^\dagger_{\rm eff}$ fit outputs — \statC{}; согласованный $g^\dagger_0$ reference — \statB{}/\statC{}; закон среды/истории $\Xi$, UDG-замыкание и экранирование конечного тела — \statO.\\
\textbf{Где это в полном препринте.} Часть~II, разделы о
галактической динамике, radial acceleration relation (RAR),
подгонках к данным SPARC и приложении по UDG-стресс-тесту.\\
\textbf{Дискриминант.} Надёжный наклон BTFR, несовместимый с $4$, отвергает принятое deep-branch замыкание.  Отдельно заранее специфицированный boundary-aware Jeans likelihood для объявленной UDG-выборки проверяет это замыкание, а не ECT целиком.  Неограниченное семейство $\Xi$ неидентифицируемо и потому не используется как фальсификатор.
\end{mechbox}

% =====================================================================
\section{Калибровочный сектор и завершение SU(3)}
\label{sec:gauge}
% =====================================================================

На гладкой когерентной ветви компактная фаза $\theta$ допускает
ковариантность расслоения, связь которой играет роль эффективного
калибровочного поля $U(1)$; целочисленные числа обмотки вокруг
дефектных ядер тогда влекут абелеву решетку зарядов компактной
группы $U(1)$. Полный наблюдаемый спектр зарядов Стандартной модели —
включая гиперзаряд, структуру ароматов и поколений — из этого
механизма не выводится и остается \statO.
Зафиксированный слой согласованности \emph{colour-interlock}
сейчас ограничивает возможные завершения: при стандартном
глобальном факторе типа $\mathbb{Z}_6$ данные по гиперзаряду,
цвету и слабому взаимодействию должны удовлетворять явному
сравнению по модулю, которое фиксирует дробный электрический
паттерн для каждого цветового класса
($Q\equiv -t/3 \pmod 1$) и сводит остаточную свободу к явно
характеризованному дискретному семейству. Это требование
согласованности для завершения, а не вывод.
Хиральный сектор $SU(2)$ структурно связан с ориентированной
нарушенной фазой через ориентационную переменную $n_A$ и служит
предметом послойной электрослабой программы
(OP-EW-scale, OP-EW-locking, OP-EW-gauge, OP-EW-naturalness,
OP-EW-matter). Условия аномалий в электрослабом секторе для
эмерджентного хирального маршрута проверены на унаследованном
содержании материи (\statA), а связь с предпочтительным
направлением не влияет на коэффициенты хиральных аномалий.

Цветовой сектор структурно отличен.
Алгебраическое препятствие $\mathfrak{su}(3)\not\subset
\mathfrak{so}(4)$ показывает, что точный локальный $SU(3)$ не
может возникнуть из одного только минимального
$O(4)$-упорядоченного конденсата. Поэтому ECT вводит как
отдельный обратно сконструированный слой:

\textbf{Постулат~P7.}
$\Phi$-среда снабжается локальным комплексным внутренним модулем
ранга~3, эрмитовой формой и фиксированной комплексной формой
объема. Соответствующая избыточность реперов равна $U(3)$;
фиксированная комплексная форма объема сокращает ее до $SU(3)$.

Далее в части~IV доказываются три obstruction-теоремы:
отсутствие древовидной массы глюона, невозможность цветового
механизма Хиггса
маршрута и тот факт, что осцилляторная башня цветового сектора не
является физическим цветовым спектром.
На уровне аномалий одни только условия сокращения аномалий не
выбирают цветовой ранг (они требуют лишь нечетного числа цветов),
поэтому вход ранга~3 в P7 должен приходить либо из
унаследованных данных о зарядах, либо из открытой программы
нулевых мод.
Установлен и центральный негативный результат: при постоянных
фоновых константах связи завершение цветового сектора P7
является структурно \emph{классифицирующим, а не
предсказательным}; оно не генерирует самостоятельной
предсказательной физики сильного сектора сверх простой
переустановки соглашений о согласовании параметров.
Топологическая связь $\Theta(\phi)F\tilde{F}$ в рамках ECT
\emph{возвращает} strong-CP проблему, а не решает ее.

\begin{mechbox}
\textbf{Механизм в ECT.}
Абелева решетка зарядов компактной группы $U(1)$ из
ковариантности расслоения компактной фазы; хиральный $SU(2)$ из
ориентированной нарушенной
фазы; точный локальный $SU(3)$ только через дополнительный
внутренне-модульный постулат P7.\\
\textbf{Что этот результат устанавливает.} Производную
решетку компактной группы $U(1)$, структурный маршрут к
хиральному $SU(2)$
и четко отделенное P7-завершение для точного локального $SU(3)$.\\
\textbf{Статус.} Решетка зарядов компактной группы $U(1)$ имеет статус
\statA/\statB{} на гладкой когерентной ветви; хиральный $SU(2)$
структурно — \statB; $SU(3)$ постулируется (P7); полный спектр
заряда/ароматов/поколений Стандартной модели, конфайнмент и
strong-CP проблема — \statO.\\
\textbf{Где это в полном препринте.} Часть~I (компактная фаза,
$U(1)$, неабелево вложение), часть~IV ($SU(3)$-завершение, P7,
теоремы-препятствия).\\
\textbf{Фальсификатор.} Подтвержденный низкоэнергетический эффект
в цветовом секторе, требующий структуры, несовместимой с
внутренне-модульной конструкцией P7 или с точным локальным
$SU(3)$ в ECT-завершении, фальсифицировал бы текущий слой
цветового сектора.
\end{mechbox}

% =====================================================================
\section{Фальсификаторы и дискриминанты}
\label{sec:falsifiers}
% =====================================================================

Препринт содержит явный реестр фальсификаторов и второй реестр
будущих дискриминантов ND1--ND13 в таблицах предсказаний и
дискриминантов. Наиболее информативные фальсификаторы и
дискриминанты текущего состояния ECT таковы:

\begin{enumerate}[leftmargin=1.4em,itemsep=2pt]
\item \textbf{F1 — исследовательская зависимость $g^\dagger$ от среды/истории.}
      Подогнанный масштаб $g^\dagger_{\rm eff}$ можно регрессировать
      по заранее объявленным ковариатам среды и истории как тест
      открытого закона~$\Xi$. Знак и амплитуда сейчас не предсказаны;
      нулевой результат отвергает только проверенный анзац $\Xi$,
      а корреляция сама по себе не выводит ECT.
\item \textbf{Скорость гравитационных волн / универсальность
      лоренцева конуса.}
      Текущая эффективная параметризация согласуется с
      ограничением GW170817
      $|c_{\rm GW}/c-1|<6\times10^{-15}$;
      планковски подавленная оценка задержки имеет вид
      $\Delta t_{\rm GW}\sim 10^{-64}\,\mathrm{s}$.
      Надежно установленная секторно-зависимая низкоэнергетическая
      скорость за пределами этой границы фальсифицировала бы
      возникновение универсального конуса.
\item \textbf{ND-13 — гранично-корректный UDG inverse test.}
      Для DF4 при центральных Wolf-proxy входах $\Xi_{\rm alg}<0$, поэтому
      внутри принятого замыкания нет физического решения $\Xi\ge0$. Полный
      Jeans likelihood проверяет это замыкание, а не ECT целиком.
\item \textbf{Gravity-induced entanglement (класс BMV).}
      Операциональный критерий отделяет режим различимости по
      $\phi$-каналу (без GIE) от ориентационного режима,
      потенциально порождающего запутывание.
\item \textbf{Гравитационный канал записи.}
      Текущие мезоскопические значения, включая $38$\,мс для
      микронной наносферы, являются импортированными бенчмарками
      Ди\'оси--Пенроуза.  Физическая вершина ECT, зависящее от
      состояния шумовое ядро, заряд конечного тела и абсолютная
      скорость остаются Open/Not Identifiable.  Старые строки
      кроссовера, помеченные только $(P,T)$, невоспроизводимы при
      заявленной ideal-gas N$_2$ конвенции.
\item \textbf{Отсутствие перехода в область $w<-1$.}
      Замыкание упорядоченной ветви дает $w(z)\ge -1$ при всех
      красных смещениях в текущем секторе; статистически надежное
      обнаружение $w<-1$ поставило бы замыкание под давление.
\item \textbf{Негативный дискриминант: сектор $S_8$ / слабое
      линзирование.}
      Текущий $\varepsilon$-механизм не рассчитан на снятие
      напряжения $S_8$; этот сектор трактуется как негативный
      дискриминант текущего эффективного слоя.
\end{enumerate}

Дополнительные дискриминанты — включая гравитационную память и
структуру мягких гравитонов, испарение черных дыр, совместимое с
кривой Пейджа, и локализованные возбуждения в дефектных ядрах —
каталогизированы в полном препринте.

ECT в нынешнем виде не заявляет независимых предсказаний новых
частиц сверх уже известных степеней свободы: любые новые возбуждения,
упомянутые выше, обсуждаются как дополнительные кандидатные
каналы возбуждения $\Phi$-среды, а не как установленные
предсказания новых частиц.

% =====================================================================
\section{Сводная карта статусов}
\label{sec:status_map}
% =====================================================================

\renewcommand{\arraystretch}{1.22}
\small
\footnotesize
\setlength{\tabcolsep}{3.0pt}
\begin{longtable}{@{}L{3.2cm} L{4.1cm} L{2.7cm} L{3.6cm}@{}}
\caption{Сводная карта статусов ECT. В колонке
``Открытая зависимость'' указан ближайший фальсификатор или
открытая зависимость, способная ослабить статус строки.
Указатели на разделы даны в соответствующих mechanism boxes;
полные технические локации приведены в препринте.}
\label{tab:status_map}\\
\toprule
\textbf{Результат} & \textbf{Механизм} & \textbf{Статус} &
\textbf{Открытая зависимость}\\
\midrule
\endfirsthead
\multicolumn{4}{@{}l}{\textit{Таблица~\ref{tab:status_map},
продолжение.}}\\
\toprule
\textbf{Результат} & \textbf{Механизм} & \textbf{Статус} &
\textbf{Открытая зависимость}\\
\midrule
\endhead
\bottomrule
\endfoot
Лоренцев конус &
SSB упорядоченной ветви; $\alpha>\beta$ &
\statA/\statB &
секторная зависимость класса GW170817\\
Когерентный квантовый сектор; $S_0\leftrightarrow\hbar$ — Level-B matching &
$S_{\rm loop,min}=2\pi S_0$; OS/RP-мост &
\statB &
независимый стандарт малой петли для действия\\
Скалярный метрический отклик / ориентационное расширение &
Скалярный анзац; внешний прокси &
\statB{}; \statO{} &
Ковариантное $S_n$; $T^{(n)}_{\mu\nu}$\\
Развязка прямого CC-канала &
Фиксированный кинетический детерминант из $n^An_A=1$ &
\statA &
спектральное условие на оператор флуктуаций\\
Галактические RAR / BTFR &
Принятая интерполяция; условная алгебра глубокой ветви &
\statC{} fit / \statB{}/\allowbreak\statC{} наклон &
Физический координатный мост; выбор ветви\\
Скалярное экранирование / закон среды &
Локальный трейс-отклик; исследовательский $\Xi$ &
\statO{} &
Заряд тела + PPN-профиль; знак/амплитуда $\Xi$\\
Бенчмарк Ди\'оси--Пенроуза / канал записи ECT &
Импортированный D--P функционал; ядро ECT не идентифицировано &
\statC{}/\allowbreak\statO{} &
Физический проектор, состояние, шум и нормировка\\
Космологическое $\varepsilon$-окно &
Однородная эффективная поздневременная деформация &
\statC &
извлечение JWST при $\varepsilon\le 0$\\
Веса типа Борна / измерение &
RP/OS-гильбертов маршрут + Глисон; C1--C3 приняты &
\statB{} условно / \statO &
вывод C1--C3; проблема единственного исхода и обновления\\
Спин--статистика (свободное когерентное ядро) &
Евклидов OS-маршрут (класс L\"uscher--OS) &
\statA{} свободное ядро &
завершение с взаимодействиями\\
Калибровочный сектор / P7 &
маршрут компактной группы $U(1)$, хиральный $SU(2)$; $SU(3)$ через P7 &
A/B для $U(1)$; B для $SU(2)$; P7 для $SU(3)$ &
полный спектр СМ (дискретное семейство, ограниченное interlock-условием); конфайнмент; strong-CP проблема\\
\end{longtable}
\normalsize

Дополнительные статусные записи для принципа запрета Паули,
наблюдаемого значения $\rho_\Lambda$, структурной нормировки
$g^\dagger_0$ и согласования низкоэнергетической квантово-полевой
архитектуры отслеживаются в полном препринте.

% =====================================================================
\section{Ресурсы и приглашение к обсуждению}
\label{sec:resources}
% =====================================================================

\textbf{Полный технический препринт.}
\emph{Euclidean Condensate Theory (ECT): Emergence of Spacetime,
Quantum Mechanics, and Gravity from Spontaneous $O(4)$ Symmetry
Breaking}, V.~Blagovidov, 2026, 784 стр. в замороженной proposal-сборке Round~58 (до релиза число страниц может измениться).\\
\href{https://doi.org/10.5281/zenodo.18917929}{doi:10.5281/zenodo.18917929}

\textbf{Сопроводительный обзор}
(концептуальный гид по полному препринту), около $104$ стр.\\
\href{https://doi.org/10.5281/zenodo.19430795}{doi:10.5281/zenodo.19430795}

\textbf{Код и скрипты генерации фигур.}\\
\href{https://github.com/chufelo/ECT-preprint-code}{github.com/chufelo/ECT-preprint-code}
(подгонки SPARC, анализ сохраняемого $\varepsilon$-окна,
UDG-стресс-тест, графики конденсатных масштабов, сравнительная
временная шкала).

\textbf{Контакты автора.}
Valeriy Blagovidov,
\texttt{blagovidov@phystech.edu},
\texttt{vblagovidov@gmail.com}.\,
ORCID:
\href{https://orcid.org/0009-0008-6707-7068}{0009-0008-6707-7068}.

\textbf{Приглашение.}
Наиболее полезный тип обратной связи для скептически настроенного
читателя этого документа — понять, не происходит ли с любым из
семи структурных утверждений из \S\ref{sec:scope} одно из трех:
(i)~оно противоречит установленному результату,
(ii)~для него существует более простой механизм в уже известной
рамке, или (iii)~оно опирается на скрытое вспомогательное
предположение, не отмеченное здесь.
Особенно полезна обратная связь по явно открытым пунктам из
\S\ref{sec:scope} — строгой теореме Борна из первых принципов и
проблеме единственного исхода, статусу Bell/CHSH/Tsirelson,
наблюдаемому значению $\rho_\Lambda$, закону среды $\Xi$,
динамическому происхождению упорядоченной ветви (OP-grad),
гиперзаряду и полной charge/flavour/generation-структуре
Стандартной модели, а также strong-CP проблеме.

% =====================================================================
\bibliographystyle{plain}
\bibliography{../references}
% =====================================================================

\end{document}
